%%
%% AgentTelemetry: An OpenTelemetry-Native Observability SDK
%% for AI Agent Systems
%%
%% Target: ICSE 2027 Tool Demonstration and Data Showcase Track
%% Page limit: 4 pages, including references.
%% Format: IEEEtran 10pt conference (no compsoc), single-anonymous.
%%

\documentclass[10pt,conference]{IEEEtran}

\usepackage[utf8]{inputenc}
\usepackage[T1]{fontenc}
\usepackage{cite}
\usepackage{booktabs}
\usepackage{array}
\usepackage{graphicx}
\usepackage{xcolor}
\usepackage{listings}
\usepackage{pifont}
\usepackage{url}
\usepackage{hyperref}
\hypersetup{colorlinks=true,linkcolor=black,citecolor=black,urlcolor=blue!60!black}

%% Allow line breaks inside \texttt{...} on underscores, dots, hyphens.
\makeatletter
\g@addto@macro\UrlBreaks{\do\_\do\.\do\-\do\,\do\?\do\/}
\makeatother
\sloppy

%% Stretch slightly so 4 pages set cleanly.
\setlength{\emergencystretch}{2em}
\tolerance=1500
\hbadness=10000

%% --- Code listing style (compact) ---------------------------------
\definecolor{lstkw}{HTML}{0033B3}
\definecolor{lststr}{HTML}{067D17}
\definecolor{lstcm}{HTML}{8C8C8C}
\lstset{
  basicstyle=\ttfamily\scriptsize,
  keywordstyle=\color{lstkw}\bfseries,
  commentstyle=\color{lstcm}\itshape,
  stringstyle=\color{lststr},
  language=Python,
  showstringspaces=false,
  columns=fullflexible,
  keepspaces=true,
  breaklines=true,
  aboveskip=2pt,
  belowskip=2pt,
  xleftmargin=0.6em,
  framexleftmargin=0pt
}

\begin{document}

\title{AgentTelemetry: An OpenTelemetry-Native\\Observability SDK for AI Agent Systems}

\author{\IEEEauthorblockN{Krishna Chaitanya Balusu}
\IEEEauthorblockA{Independent Researcher\\
San Francisco, USA\\
krishnabkc15@gmail.com}}

\maketitle

\begin{abstract}
LLM-based AI agents fail in modes---circular delegation, planning
loops, silent guardrail bypass, stale memory---that existing
observability tools cannot represent. OpenTelemetry's GenAI semantic
conventions cover LLM invocation and tool execution but leave five
agent orchestration phases (planning, reasoning, safety,
delegation, memory) without span-level vocabulary, so the resulting
traces show \emph{what} the model said but not \emph{what the agent did}.
We present \textsc{AgentTelemetry}, an Apache-licensed,
pip-installable Python SDK that adds nine agent-specific span
kinds, seven framework auto-instrumentors (LangChain, CrewAI,
AutoGen, OpenAI, Anthropic, LlamaIndex, plus a manual API), three
privacy levels, and four analysis modules
(anomaly detection, cost aggregation, decision attribution,
hallucination tracing). The SDK is OpenTelemetry-native and exports
to any OTLP backend (Jaeger, Grafana Tempo, Datadog, etc.). A
developer instruments an existing LangChain agent in five lines and
immediately gains structural visibility into coordination failures.
A companion 3{,}780-row public benchmark (Zenodo
DOI~10.5281/zenodo.20129005) shows that the nine-kind span
vocabulary can express fault classes that vanilla OpenTelemetry
cannot. All artifacts---PyPI release, source, dataset, and a
3--5~minute screencast---are openly available.
\end{abstract}

\begin{IEEEkeywords}
AI agents; observability; OpenTelemetry; fault detection;
distributed tracing; LLM monitoring; software tools
\end{IEEEkeywords}

\section{Introduction}
\label{sec:intro}

Autonomous LLM agents now plan tasks, call tools, delegate to
sub-agents, retrieve from vector stores, and reason over multi-turn
state. When they fail, however, the trace a developer sees in their
observability backend rarely reflects the agent's actual cognitive
trajectory. OpenTelemetry's GenAI semantic conventions~\cite{otel,
otel_genai} model only LLM invocation and tool execution. Failures
that originate from \emph{coordination} (circular delegation,
agent misroute, infinite reasoning loops) or from \emph{cognitive
state} (stale retrieval, memory corruption, guardrail
bypass)~\cite{mast, agentops_paper} therefore appear as
indistinguishable nests of \texttt{gen\_ai.completion} spans.

Existing commercial and open-source agent observability
tools---LangSmith~\cite{langsmith}, Langfuse~\cite{langfuse},
AgentOps~\cite{agentops}, Arize Phoenix~\cite{phoenix},
OpenLIT~\cite{openlit}---either reuse the generic GenAI spans or
introduce closed proprietary schemas (\cite{openinference}
partially excepted). None offers a typed span vocabulary for the
five missing orchestration phases, and several are SaaS-only.

\textsc{AgentTelemetry} closes this gap with an OpenTelemetry-native
SDK. A developer who already understands OTel does not need to
learn a new dashboarding product or accept a hosted backend; they
\texttt{pip install agenttelemetry}, add two lines, and their
existing Jaeger or Tempo instance starts showing typed spans for
\texttt{PLANNING}, \texttt{REASONING}, \texttt{DELEGATION},
\texttt{GUARD\_RAIL}, and \texttt{MEMORY}. The same span graph then
feeds four offline analysis modules that flag known agent failure
patterns structurally, without inspecting prompt content.

\noindent\textbf{Contributions.} This tool demonstration paper
contributes (i)~a pip-installable, Apache-2.0 SDK
(4{,}100~LOC, 111 tests) implementing nine agent span kinds and
seven framework adapters; (ii)~four analysis modules that
operate over the resulting traces; (iii)~a publicly downloadable
3{,}780-row fault detection benchmark on Zenodo that demonstrates
the SDK produces a detectable signal across 14~fault classes; and
(iv)~a 3--5~minute screencast walking through end-to-end
instrumentation, trace inspection in Jaeger, and anomaly detection.

\section{Tool Overview}
\label{sec:overview}

\begin{figure}[t]
\centering
\setlength{\fboxsep}{0pt}\setlength{\fboxrule}{0.4pt}
\fbox{%
\begin{minipage}{0.96\columnwidth}
\centering\vspace{2pt}\ttfamily\scriptsize
\begin{tabular}{c}
\textbf{User Code (LangChain, CrewAI, AutoGen, custom)}\\
$\downarrow$ \rmfamily\itshape callback / monkey-patch / context mgr\\
\fbox{\rmfamily\bfseries\small Adapter (7 framework instrumentors)}\\
$\downarrow$\\
\fbox{\rmfamily\bfseries\small AgentTelemetry SDK}\\
\fbox{\small spans \,$|$\, privacy \,$|$\, exporter}\\
$\downarrow$ \rmfamily\itshape OTLP/HTTP \,$|$\, OTLP/gRPC \,$|$\, console \,$|$\, JSON\\
\fbox{\rmfamily\bfseries\small Backend: Jaeger, Tempo, Datadog, OTLP-any}\\
$\downarrow$ \rmfamily\itshape trace export\\
\fbox{\rmfamily\bfseries\small Analysis: anomaly \,$|$\, cost \,$|$\, attribution \,$|$\, halluc.}
\end{tabular}\vspace{2pt}
\end{minipage}}
\caption{AgentTelemetry architecture. User code is instrumented
through a framework adapter or the manual API. Spans flow through
the SDK's privacy filter and exporter to any OTLP-compatible
backend. Analysis modules consume completed traces offline to
flag structural failure patterns.}
\label{fig:arch}
\end{figure}

\textbf{Span kinds.} The SDK defines an
\texttt{AgentSpanKind} enum with nine values:
\texttt{AGENT}, \texttt{LLM\_CALL}, \texttt{TOOL\_CALL},
\texttt{PLANNING}, \texttt{REASONING}, \texttt{RETRIEVAL},
\texttt{GUARD\_RAIL}, \texttt{DELEGATION}, and \texttt{MEMORY}.
Each is exported as the standard
\texttt{agent.span.kind} OTel attribute, so existing OTLP backends
display them without modification. Three of the nine
(\texttt{LLM\_CALL}, \texttt{TOOL\_CALL}, \texttt{RETRIEVAL})
overlap conceptually with OTel GenAI; the remaining six fill the
orchestration gap.

\textbf{Adapters.} Table~\ref{tab:adapters} summarises the seven
framework adapters. Each adapter uses lazy imports: the framework
package is required only at \texttt{instrument()} time, so users do
not pay an import cost for frameworks they don't use.

\begin{table}[t]
\caption{Framework adapters shipped with AgentTelemetry 0.1.0.}
\label{tab:adapters}
\centering\footnotesize
\setlength{\tabcolsep}{4pt}
\begin{tabular}{@{}l l l r@{}}
\toprule
\textbf{Framework} & \textbf{Class} & \textbf{Strategy} & \textbf{LOC}\\
\midrule
LangChain     & \texttt{LangChainInstr.}  & Callback handler & 359\\
CrewAI        & \texttt{CrewAIInstr.}     & Hook system      & 264\\
AutoGen       & \texttt{AutoGenInstr.}    & Monkey-patch     & 256\\
Anthropic SDK & \texttt{AnthropicInstr.}  & Monkey-patch     & 294\\
OpenAI SDK    & \texttt{OpenAIInstr.}     & Monkey-patch     & 245\\
LlamaIndex    & \texttt{LlamaIndexInstr.} & Span handler     & 280\\
Custom (manual) & \texttt{start\_agent\_span} & Context mgr  & 315\\
\bottomrule
\end{tabular}
\end{table}

\textbf{Privacy.} Three levels (\texttt{NONE},
\texttt{METADATA\_ONLY} [default], \texttt{FULL}) gate what
attribute payloads reach the exporter; metadata-only retains
structural diagnostics (model names, token counts, USD cost,
status) while suppressing prompt and completion text.

\textbf{Analysis modules.} Four modules consume completed traces
offline:
\texttt{AnomalyDetector} flags circular delegation, infinite
retries, cost explosions, and context overflow by walking the
span graph;
\texttt{CostAggregator} sums per-model and per-agent USD spend;
\texttt{DecisionAttributor} links each tool call back to the LLM
decision that triggered it;
\texttt{HallucinationTracer} marks LLM outputs whose tokens are
not grounded in spans of kind \texttt{RETRIEVAL} immediately
preceding them. All four operate on serialized
trace JSON, so they can run on Jaeger exports, persisted
\texttt{.jsonl} bundles, or live OTel collector pipelines.

\section{Usage}
\label{sec:usage}

Listing~\ref{lst:manual} shows the entire developer-facing surface
for a hand-instrumented agent: one \texttt{configure} call and a
nested set of context managers. The same trace can be exported to
Jaeger by setting \texttt{OTEL\_EXPORTER\_OTLP\_ENDPOINT}; no
SDK-side change is required.

\begin{lstlisting}[caption={Manual instrumentation in 10 lines.},label={lst:manual}]
from agenttelemetry import (
    configure, start_agent_span, AgentSpanKind)

provider = configure(service_name="my-agent", console=True)
tracer = provider.get_tracer()

with start_agent_span("research", AgentSpanKind.AGENT, tracer=tracer):
    with start_agent_span("call-gpt4", AgentSpanKind.LLM_CALL,
                          tracer=tracer) as llm:
        llm.set_attribute("llm.model", "gpt-4o")
        llm.set_attribute("llm.input_tokens", 500)
    with start_agent_span("search", AgentSpanKind.TOOL_CALL,
                          tracer=tracer) as tool:
        tool.set_attribute("tool.name", "google")
\end{lstlisting}

Listing~\ref{lst:auto} shows the equivalent for an existing
LangChain~\cite{langchain} ReAct~\cite{react} agent. The
\texttt{LangChainInstrumentor} attaches a callback handler that
maps LangChain's internal events (chain start, LLM start, tool
start, agent action) onto the correct \texttt{AgentSpanKind}, so
\emph{no application code changes}.

\begin{lstlisting}[caption={Auto-instrumentation for an existing LangChain agent.},label={lst:auto}]
from agenttelemetry import configure
from agenttelemetry.adapters import (
    LangChainInstrumentor)

provider = configure(service_name="my-agent")
LangChainInstrumentor().instrument(
    tracer_provider=provider.tracer_provider)
# existing LangChain code runs unchanged
\end{lstlisting}

A trace produced by Listing~\ref{lst:auto} on a tool-using ReAct
agent appears in Jaeger as the nested span graph
\texttt{AGENT} $\to$ \texttt{PLANNING} $\to$ \texttt{LLM\_CALL}
$\to$ \texttt{TOOL\_CALL} $\to$ \texttt{REASONING} $\to$
\texttt{LLM\_CALL}, with the orchestration spans
appearing alongside the standard GenAI spans rather than replacing
them. This is the key visual payoff of the SDK and is the
centerpiece of the accompanying screencast (Sec.~\ref{sec:avail}).

\textbf{Walk-through: catching a circular delegation.}
Consider a CrewAI swarm (\path{examples/multi_agent.py} in the
repository) in which agent~A delegates a
sub-task to agent~B, which routes back to~A on an unrecognized
intent. Without typed spans the resulting trace is a deep nest of
LLM completions and tool calls---no single span attribute
signals the cycle. Auto-instrumented with \texttt{CrewAIInstrumentor},
the same run produces the span graph shown in
Listing~\ref{lst:trace}: every cross-agent hand-off emits a
\texttt{DELEGATION} span carrying
\texttt{agent.delegation.target} and \texttt{agent.delegation.depth}
attributes. The bundled \texttt{AnomalyDetector} traverses this
graph with three lines of code and returns a list of typed
\texttt{Anomaly} records (Listing~\ref{lst:trace}, bottom). The
same call also flags cost explosions and infinite retries from the
same trace. This end-to-end loop---instrument, export,
analyse---runs in under 30~seconds on a developer laptop and is
the centre of the demo video.

\begin{lstlisting}[caption={Jaeger-style trace dump and \texttt{AnomalyDetector} output for the circular-delegation case.},label={lst:trace}]
AGENT  researcher           [span-a1]
 PLANNING  decide-route     [span-a2]
  LLM_CALL  gpt-4o          [span-a3]
 DELEGATION  -> writer      [span-a4] target=writer depth=1
  AGENT  writer             [span-b1]
   DELEGATION  -> researcher [span-b2] target=researcher depth=2
    AGENT  researcher       [span-a1]   <-- cycle

>>> AnomalyDetector().detect(trace)
[Anomaly(type='CIRCULAR_DELEGATION',
         evidence=['span-a1','span-b2','span-a1'],
         depth=3, severity='high')]
\end{lstlisting}

\section{Comparison with Existing Tools}
\label{sec:compare}

Table~\ref{tab:compare} contrasts AgentTelemetry with the five
most prominent agent-observability tools as of mid-2026. The
per-competitor wedge is the load-bearing claim of this section:
against \textbf{LangSmith}~\cite{langsmith} (closed-source SaaS,
no OTLP egress), AgentTelemetry is self-hostable Apache code that
exports to any OTLP backend the team already runs.
Against \textbf{Langfuse}~\cite{langfuse} (OTel-native but generic
LLM spans), AgentTelemetry adds five typed orchestration span
kinds (\texttt{PLANNING}, \texttt{REASONING}, \texttt{DELEGATION},
\texttt{GUARD\_RAIL}, \texttt{MEMORY}) plus an \texttt{AGENT} root
span that Langfuse currently folds into untyped trace events.
Against \textbf{AgentOps}~\cite{agentops} (OSS dashboard product),
AgentTelemetry ships an offline analysis library
(\texttt{AnomalyDetector}, \texttt{CostAggregator},
\texttt{DecisionAttributor}, \texttt{HallucinationTracer}) that
operates on exported trace JSON without a hosted backend.
Against \textbf{Phoenix}~\cite{phoenix} (OpenInference-based),
AgentTelemetry stays on the OpenTelemetry semantic-conventions
track and contributes a draft OTEP
(\path{otep/agent-semantic-conventions.md}) rather than maintaining
a parallel schema. Against the closest competitor,
\textbf{OpenLIT}~\cite{openlit} (Apache, OTel-native, 20+
provider instrumentations), the wedge is structural: OpenLIT
instruments individual \emph{provider SDKs} (Anthropic, OpenAI,
vector stores) and emits standard \texttt{gen\_ai.*} spans, so a
CrewAI delegation cycle is invisible at the span-kind level.
AgentTelemetry instruments \emph{agent frameworks} (LangChain,
CrewAI, AutoGen) and emits typed orchestration spans, so the same
cycle shows up directly as a \texttt{DELEGATION} subgraph that
\texttt{AnomalyDetector} can match without parsing prompt text.

\begin{table*}[t]
\caption{AgentTelemetry vs.\ existing agent-observability tools (verified from each project's official documentation, May 2026). \ding{51}~=~yes, \ding{55}~=~no.}
\label{tab:compare}
\centering\footnotesize
\setlength{\tabcolsep}{3.5pt}
\begin{tabular}{@{}l c c l c l c c@{}}
\toprule
\textbf{Tool} & \textbf{OSS} & \textbf{OTel} & \textbf{Typed agent spans} & \textbf{Adapters} & \textbf{Analysis} & \textbf{Open semconv} & \textbf{Self-host}\\
 & & \textbf{native} & & & & \textbf{proposal} & \\
\midrule
LangSmith~\cite{langsmith}   & \ding{55} (SaaS)    & \ding{55}                & generic LLM/tool   & many   & proprietary            & \ding{55} & \ding{55}\\
Langfuse~\cite{langfuse}     & \ding{51}           & \ding{51}                & generic LLM        & 50+    & dashboards             & \ding{55} & \ding{51}\\
AgentOps~\cite{agentops}     & \ding{51}           & not docu.                & generic            & many   & dashboards             & \ding{55} & \ding{51}\\
Phoenix~\cite{phoenix}       & \ding{51}           & via OpenInf.             & generic LLM/tool   & many   & evaluation             & \ding{55} (OpenInf., not OTel) & \ding{51}\\
OpenLIT~\cite{openlit}       & \ding{51}           & \ding{51}                & generic LLM        & 20+    & dashboards             & \ding{55} & \ding{51}\\
\textbf{AgentTelemetry (this work)} & \ding{51} (Apache) & \ding{51}     & \textbf{9 typed kinds} & 7      & \textbf{4 modules}     & \ding{51} (OTEP draft) & \ding{51}\\
\bottomrule
\end{tabular}
\end{table*}

\section{Evidence of Detectable Signal}
\label{sec:eval}

A 3{,}780-row controlled benchmark accompanies the SDK release
(3{,}528 fault-injection runs $+$ 252 fault-free
controls)~\cite{agenttelemetry_aiware, agenttelemetry_zenodo} and
is publicly downloadable as a tab-separated file from the Zenodo
record. The benchmark covers 14 fault classes across six telemetry
conditions and seven agent frameworks. Aggregate Fault Detection
Rate (FDR), reproducible verbatim from
DOI~10.5281/zenodo.20129006:

\begin{center}\footnotesize
\setlength{\tabcolsep}{4pt}
\begin{tabular}{@{}l r r@{}}
\toprule
\textbf{Telemetry condition} & \textbf{FDR} & \textbf{Detected}\\
\midrule
No telemetry                                & 0.000 & 0/14\\
Vanilla OpenTelemetry                       & 0.429 & 6/14\\
OTel GenAI conventions                      & 0.429 & 6/14\\
OpenInference                               & 0.429 & 6/14\\
AgentTelemetry (metadata privacy)           & 0.612 & 7/14\\
AgentTelemetry (full vocabulary, upper bd.) & \textbf{1.000} & \textbf{14/14}\\
\bottomrule
\end{tabular}
\end{center}

The upper bound is the relevant claim for a tool-demonstration
audience: the nine-kind span vocabulary is structurally
\emph{sufficient} to express all 14 fault classes, including the
eight coordination, cognitive, and policy faults
(\path{circular_delegation}, \path{hallucination},
\path{stale_retrieval}, \path{guardrail_bypass},
\path{planning_failure}, \path{reasoning_loop},
\path{agent_misroute}, \path{memory_corruption}) that no
baseline can express. The 0.612 mid-row is a per-framework
conformance gap, not a metamodel limit: only the bundled
\texttt{custom\_agent} reference adapter currently emits the full
nine-kind vocabulary, while third-party adapters emit 3--7 typed
kinds out of 9. Closing this gap is ongoing work and is tracked in
the repository roadmap. Full methodology, ablations, and
statistical analysis are reported separately~\cite{agenttelemetry_aiware};
we reference the artifact here strictly as evidence that the SDK
produces a usable signal.

\textbf{Overhead.} The in-tree benchmark
\path{tests/benchmarks/test_overhead.py} (run against an
\texttt{InMemorySpanExporter}) asserts that a single instrumented
span completes in $<$5~ms and a nested
agent$\to$LLM$\to$tool triple in $<$10~ms on a 2024-class laptop
CPU, well below typical LLM round-trip latencies.

\section{Availability and Reproducibility}
\label{sec:avail}

\textbf{Install.} \texttt{pip install agenttelemetry==0.1.0}
(PyPI, Apache-2.0)~\cite{agenttelemetry_pypi}. The unpinned
\texttt{pip install agenttelemetry} also resolves to 0.1.0 today.
On a stock 2024-class laptop the install completes in
$<$30~seconds (the SDK has four runtime dependencies:
\texttt{opentelemetry-api}, \texttt{opentelemetry-sdk},
\texttt{opentelemetry-exporter-otlp}, \texttt{pydantic}). Optional
extras (\texttt{[langchain]}, \texttt{[crewai]}, \texttt{[autogen]},
\texttt{[anthropic]}, \texttt{[openai]}, \texttt{[llamaindex]},
\texttt{[all]}) pull in the corresponding framework dependency
groups.

\textbf{Source.} \url{https://github.com/Krishnachaitanyakc/AgentTelemetry}
(Apache-2.0, 4{,}100 SDK LOC, 111 tests, CI on Python 3.9--3.12).

\textbf{Archived artifact.} Zenodo concept
DOI~\href{https://doi.org/10.5281/zenodo.20129005}{10.5281/zenodo.20129005};
pinned version DOI~\href{https://doi.org/10.5281/zenodo.20129006}{10.5281/zenodo.20129006}
contains the 0.1.0 source plus the 3{,}780-row benchmark TSV and a
\texttt{requirements.lock} that reproduces every number above
verbatim.

\textbf{Screencast.} A 3--5~minute video walks through
installation, manual and auto-instrumentation, the Jaeger trace
view, and the anomaly detector on a circular-delegation example:
\url{https://youtu.be/AGENTTELEMETRY-ICSE27-DEMO}.

\textbf{Semantic conventions.} A draft OTEP proposing the nine
agent span kinds for upstream adoption lives at
\path{otep/agent-semantic-conventions.md} in the repository and
is being discussed with the OpenTelemetry Semantic Conventions
working group.

\section{Limitations and Roadmap}
\label{sec:limits}

The SDK is at \textbf{alpha} maturity (PyPI development status~3).
Three limitations bound current use. First, the 0.612 average FDR
in Sec.~\ref{sec:eval} reflects that only the bundled
\texttt{custom\_agent} reference adapter emits the full nine-kind
vocabulary; third-party framework adapters emit 3--7 of the nine
typed kinds. Conformance work for LangChain, CrewAI, and AutoGen
adapters is the top roadmap item. Second, the
\texttt{HallucinationTracer} relies on a token-overlap
heuristic against \texttt{RETRIEVAL}-kind spans and is therefore
recall-biased; semantic-similarity grounding is a planned
replacement. Third, the OTEP that proposes the nine span kinds
for upstream OpenTelemetry adoption is under community review and
not yet a ratified standard, so backends that match on
\texttt{agent.span.kind} are configuring against a draft attribute.

\section{Conclusion}
\label{sec:conclusion}

\textsc{AgentTelemetry} ships today as a one-command-installable
Python SDK that gives any OpenTelemetry-equipped team typed
visibility into the orchestration phases that distinguish AI
agents from single LLM calls. The accompanying public benchmark
and screencast let reviewers and practitioners verify, in
minutes, that the span vocabulary captures coordination and
cognitive failure modes other tools elide. We invite the
community to install the package, build new framework adapters,
file fault patterns, and---via the OTEP draft---help shape an
upstream OpenTelemetry semantic convention for agent
observability.

\bibliographystyle{IEEEtran}
\bibliography{refs}

\end{document}
